\chapter{大语言模型}
\label{chap:nlp-large-language-models}

同一份活动通知可以用来补全句子，也可以作为回答时间、地点问题的材料。若为每个问题分别建立一组词频表，就很难把一种表达中学到的规律用到另一种表达中。上一章已经说明，连续表示与参数共享能够缓解这种稀疏性：模型根据上下文计算表示，再把表示映射为语言单位的概率。本章接着回答，当语料、模型和学习预算同时扩大时，怎样构造并学习这个概率函数。

\term{大语言模型}（large language model, LLM）通常指在大规模语言材料上训练、能够将同一套参数用于多种语言活动的模型；“大”没有适用于所有时期和结构的固定参数阈值。这里以基于Transformer的语言模型为主线。所谓\term{预训练}（pre-training），是先利用广泛材料建立可复用的模型参数，再将这些参数作为后续任务学习的起点。文本本身能够提供预测目标，但材料中出现什么、模型能够看见什么、哪些预测被计分，仍然需要人为规定。

我们沿用“读书会改到周五下午两点举行”这一教学短句。模型既要从“改到周五”预测接下来的时间表达，也要利用前面的活动名称保持主题一致。要解释这种学习怎样发生，必须把四件事连起来：材料处理决定训练分布；网络结构决定从上下文到概率的计算；学习目标决定哪些预测被计分；优化过程决定有限预算内参数怎样变化。这四者共同产生预训练后的参数，任何一项的改进都不能独立保证模型会遵守用户指令。

\section{预训练语料与样本准备}
\label{sec:nlp-llm-data}

概率模型拟合的是训练过程中实际看到的序列分布。网页数量很多，不等于模型得到的语言、事实和推理材料就均衡；同一篇通知被转载一百次，也没有增加一百种独立的事件。准备语料的任务，是从原始材料中形成有边界、可追踪且能够按指定比例取用的训练样本。模型计算尚未开始，学习对象已经在这些选择中发生变化。

\subsection{来源选择与能力覆盖}
\label{sec:nlp-llm-data-sources}

设候选材料来自网页、书籍、代码和专业文档等来源。网页能够包含日常表达和近期信息，代码能够包含形式化约束与程序结构，专业文档能够包含通用网页较少涉及的概念。这里的“来源”描述材料怎样产生，并不自动给出质量或能力标签：网页可能是完整教程，也可能只有导航菜单；代码中也可能有过时接口或错误示例。来源标签的作用，是让后续处理能够检查某种材料是否被过度删除或过度重复。

为了把“覆盖”说清楚，可以为每篇文档保留语言、时间、来源、主题和长度等属性，再分别统计它们的分布。若希望模型处理中文活动通知和英文技术文档，至少需要检查两类材料是否存在、包含哪些表达以及是否覆盖相关时间范围。材料中有某个主题，仍不意味着模型能正确完成该主题的任务，因为学习还取决于采样次数、目标和参数容量。覆盖检查排除的是明显的信息缺口，不是能力保证。

材料的时间边界尤其影响事实性判断。假设训练语料中只有“读书会于周四举行”，而新的更改通知没有进入语料，降低旧材料上的预测损失不能使模型知道周五的新安排。这类缺口需要新材料或者使用时提供外部证据，不能只通过增大网络消除。为比较不同模型，应记录其材料时间范围，避免把接触过更新信息误判为相同信息条件下的建模优势。

来源选择还需要确定哪些材料可以进入当前训练集，以及敏感内容和个人信息如何处理。这些约束改变可使用的文档集合，因而属于数据准备本身。保留来源记录有助于发现问题后追溯处理规则；后续安全示范与偏好学习另有职责，不能替代这里的材料检查。

\subsection{质量筛选与重复控制}
\label{sec:nlp-llm-data-filtering}

候选文档中常混有乱码、模板、极短片段和重复正文。若直接按文档大小取样，模型会花费大量更新去拟合这些重复或低信息量的模式。质量筛选先依据明确规则或评分保留材料；重复控制再限制相同内容在训练分布中的权重。两项操作都改变训练对象，不能仅被看成不影响学习的文件清理。

规则筛选可以检查文本可解析性、语言比例、重复行或异常字符。模型评分则先定义希望保留的特征，例如文档是否连贯、是否包含解释性内容，再用分类器或语言模型估计分数。评分器的判断也有误差：只偏好规范书面语，可能删除口语、方言和少见专业表达。阈值应通过被保留与被删除样本的抽查来解释，不能把较高分数直接当作真实知识更正确。

完全重复可以通过规范化后的文本哈希识别。近似重复则需要比较文档或片段的重合，例如把文档表示为连续词元片段集合，再依据集合重合率判定是否属于同一重复簇。去重粒度会改变结果：只删完全相同的整篇文档，会留下含不同页眉的转载；激进删除共享段落，又可能误伤不同文档都需要的定义或代码模板。训练集与评价集之间也要检查重合，但这种检查不能证明没有语义改写后的污染。

DataComp-LM将数据筛选放在共同的模型、预算和评价流程下比较，FineWeb则公开分析了过滤与去重选择。它们提供的认识是，数据处理应通过可比较的学习结果检验，而不能只报告删掉多少文本；其具体收益仍受语言范围、模型规模和评价任务约束\scite{nlp-llm-dclm2024,nlp-llm-fineweb2024}。在本章的活动通知例子中，保留一份正确正文、删除九份相同转载，会降低该通知的重复曝光，但不会自动增加模型对其他活动安排的认识。

\subsection{词元化与文档序列组织}
\label{sec:nlp-llm-tokenization}

筛选后的字符串还不是网络使用的整数序列。\term{词元}（token）是模型一次读取或预测的离散单位；\term{词元化}（tokenization）把文本转换为这些单位的序列。词元可以对应字、常见子词或字节片段。这里需要区分文字、语言学意义上的词和模型单位：“下午两点”可以被分成多个词元，也可以部分合并，具体取决于词元化规则。

设编码函数将文档$d$映射为$x_{1:T}=(x_1,\ldots,x_T)$，其中$x_t\in\mathcal V$，$\mathcal V$是有限词表。输入嵌入和输出分布都以这个词表为接口。采用子词单位时，常见片段可以共用一个编号，稀有字符串再拆成较短单位；采用字节单位时，基本字节集合固定，但同一段文本往往需要更多步。一个词元覆盖多少字符，影响序列长度和可容纳的材料，因而不能直接用不同词元化器下的每词元损失比较文本建模质量。

文档组织还要说明模型在哪里开始、在哪里结束，以及相邻文档能否彼此读取。可以在通知结尾加入结束标记，将多个短文档装入一个固定长度序列；也可以为各文档设置相互隔离的注意掩码。前者若允许跨文档读取，就改变了条件信息：后一个文档的预测能够利用前一个文档。结束标记表达边界，但本身不等于强制隔离。切分长文档时，也应检查需要跨段的信息是否仍能共同出现。

例如，假设“读书会／改到／周五／下午／两点／举行”是一次教学词元化，最大长度只有四。若永远切成前四个与后两个词元，“两点”所在样本可能失去“改到周五”的条件。重叠切分可以保留部分上下文，但会增加某些词元的重复计分。具体监督位置由\ref{sec:nlp-llm-objective-positions}小节确定；在这里先保留切分与文档边界，才能知道后续目标用了什么信息。

\subsection{多源混合与采样配比}
\label{sec:nlp-llm-data-mixture}

完成统一编码后，材料的原始数量仍不等于训练时的出现比例。设处理后的$K$类来源分别给出样本分布$P_1,\ldots,P_K$，训练时先以概率$\alpha_k$选择来源$k$，再从该来源取样。实际训练分布为
\begin{equation}
 P_{\bm{\alpha}}(x)=\sum_{k=1}^{K}\alpha_k P_k(x),
 \qquad \alpha_k\geq 0,\quad\sum_{k=1}^{K}\alpha_k=1.
 \label{eq:nlp-llm-mixture}
\end{equation}
权重$\alpha_k$描述训练曝光，而$P_k$包含该来源内部的筛选、切分和取样方式。若按词元预算分配混合，权重应按所消耗的词元计量；若按样本抽取而样本长度不同，最终词元比例通常不会等于样本比例。

假设通用文本有九十万词元，专业文档有十万词元。按原始数量抽取，专业材料约占一成；若训练词元各分一半，专业材料会被重复使用更多次。提高专业权重既增加相关模式的学习机会，也减少同一预算内通用材料的曝光，并可能增加小语料的记忆。这里没有仅由材料数量决定的最优比例，需要说明希望改善哪些评价对象。

若为每种候选比例都完整训练大模型，比较成本很高。DoReMi先在较小代理模型上，根据各来源相对于参考模型的额外损失学习重加权，再将得到的比例用于较大模型；参考损失有助于避免单纯追逐本来就更难预测的来源。RegMix改为在许多比例上训练小模型，拟合“比例到评价结果”的回归关系，再选择预测较好的比例\scite{nlp-llm-doremi2023,nlp-llm-regmix2024}。前者利用代理训练中的损失反馈，后者利用跨候选配置的预测，两者都需要检验小规模判断能否迁移到目标训练。

混合比例也可以随阶段变化，例如在较广的材料上建立通用参数后，再提高某个领域的占比。只要比例或来源发生变化，就应重新说明训练分布，而不能把整个过程概括成使用了某个固定语料集。现在我们得到了可取样的词元序列，下一节讨论模型怎样把这些序列变成条件概率。

\section{模型结构与前向计算}
\label{sec:nlp-llm-computation}

给定词元序列以后，模型需要在每个位置计算可用于预测的上下文表示。我们以仅由解码器层组成的因果Transformer为基线：预测某个后续词元时，当前位置只读取已经提供的前缀；训练时也遵守同样的可见范围。网络的共同结构是反复交替进行上下文聚合与逐位置变换，再把最终表示映射为概率。下列结构变体分别改变这个计算中的某个部件，不能因为都用于长序列就视为同一种方法。

\subsection{输入嵌入与位置表示}
\label{sec:nlp-llm-embedding}

离散编号只表示词元身份，编号大小本身没有语义距离。设词表大小为$V=|\mathcal V|$，模型宽度为$d$，可学习嵌入表为$\bm{E}\in\mathbb R^{V\times d}$。对位置$t$的词元$x_t$，取出对应行并写成列向量$\bm e_t=\bm E_{x_t,:}^{\top}\in\mathbb R^d$。相同词元在不同位置读取同一行参数，因此一处样本对该行的更新能够影响其他样本中的同一词元。它提供的是共享起点，后续上下文计算才使同一词元在不同句子中得到不同表示。

嵌入向量本身没有标明词元的相对位置。把“周四改到周五”与“周五改到周四”视为无序集合，会丢失安排变化的方向。一个做法是在输入加入位置向量；另一个做法是保留位置索引，在计算位置间的相关性时使用它。本章主要采用后一条路线中的旋转位置编码，因此初始表示可取$\bm h_t^{(0)}=\bm e_t$，位置信息在查询与键形成以后进入计算。二者不是同一运算，不能把所有位置方案都写成在嵌入上相加。

词元接口也可以与主要计算单位分开。字节潜在Transformer（Byte Latent Transformer, BLT）从字节输入出发，由局部网络编码和恢复字节，主要网络则处理根据预测不确定性划分的动态片段。可预测部分可以形成较长片段，不确定部分采用较短片段，改变的是昂贵计算分配到哪些位置\scite{nlp-llm-blt2024}。它不是在固定子词词表上换一张嵌入表；片段编码、主要网络和字节输出共同构成完整概率模型。比较此类结构时，需要同时给出原始字节数、片段数和计算预算。

\subsection{上下文聚合与键值状态组织}
\label{sec:nlp-llm-attention}

输入表示需要读到与当前预测有关的历史。以一层中的一个注意力头为例，暂省略层号。将位置表示$\bm h_t\in\mathbb R^d$分别映射成查询、键和值：$\bm q_t=\bm W_Q\bm h_t$、$\bm k_t=\bm W_K\bm h_t$、$\bm v_t=\bm W_V\bm h_t$。查询表达当前位置需要寻找什么，键用于衡量历史位置与这个需要是否匹配，值承载被汇总的信息。设查询与键维度为$d_h$，值维度为$d_v$，相应地$\bm W_Q,\bm W_K\in\mathbb R^{d_h\times d}$、$\bm W_V\in\mathbb R^{d_v\times d}$。则\term{因果注意力}（causal attention）的聚合为
\begin{equation}
 a_{tj}=\frac{\exp(\bm q_t^{\top}\bm k_j/\sqrt{d_h})}
 {\sum_{i=1}^{t}\exp(\bm q_t^{\top}\bm k_i/\sqrt{d_h})},
 \qquad
 \bm o_t=\sum_{j=1}^{t}a_{tj}\bm v_j.
 \label{eq:nlp-llm-causal-attention}
\end{equation}
权重只在$j\leq t$上归一化，保证$\bm o_t$不使用未来词元。缩放因子控制点积随维度变化的尺度，但不能单独保证注意力分布合适；投影参数还需要通过学习确定。

为了在匹配中使用相对距离，\term{旋转位置编码}（rotary position embedding, RoPE）对查询与键施加由位置决定的旋转。先考虑二维情形，取旋转角频率$\omega$，令
\begin{equation}
 \bm R_t=\begin{pmatrix}
 \cos(t\omega)&-\sin(t\omega)\\
 \sin(t\omega)&\cos(t\omega)
 \end{pmatrix},\qquad
 (\bm R_t\bm q_t)^{\top}(\bm R_j\bm k_j)
 =\bm q_t^{\top}\bm R_{j-t}\bm k_j.
 \label{eq:nlp-llm-rope}
\end{equation}
最后一步来自旋转矩阵的复合关系$\bm R_t^{\top}\bm R_j=\bm R_{j-t}$。高维实现将多个不同频率的二维旋转组成分块矩阵。匹配分数因此能够依赖相对位移$j-t$，同时仍受内容向量影响。RoPE改变的是查询与键，不要求给值向量或输入嵌入施加同样旋转；相对位置进入公式也不保证模型能直接处理任意未训练长度\scite{su2021roformer}。

\term{多头注意力}（multi-head attention, MHA）使用$H$组投影，各自计算式\eqref{eq:nlp-llm-causal-attention}，再拼接各头输出并投影回$d$维。不同头可以学习不同的内容匹配方式，但不能预先断言某个头必然负责语法、另一个头必然负责事实。聚合后还可以加入依赖当前表示的门控，例如$\tilde{\bm o}_t=\sigma(\bm W_g\bm h_t)\odot\bm o_t$，其中$\sigma$是逐坐标的sigmoid函数，$\odot$表示逐坐标乘法。Gated Attention研究了这类聚合输出门控：它改变信息经过聚合后怎样传出，通常没有因此减少需要读取的历史位置\scite{qiu2025gated}。

当模型逐步产生新词元时，旧位置的键和值可以保存，供新查询读取。若每头键和值均为$d_h$维，每层每个历史词元在MHA下需要保存$2Hd_h$个数。\term{多查询注意力}（multi-query attention, MQA）让所有查询头共享一组键和值，变为$2d_h$个数；\term{分组查询注意力}（grouped-query attention, GQA）将查询头分成$G$组，每组共享键值，变为$2Gd_h$个数。这种共享减少状态量，也约束了不同查询头可以使用的键值表示；改变既有模型的共享方式通常需要相应训练或适应\scite{shazeer2019mqa,ainslie2023gqa}。

多头潜在注意力（multi-head latent attention, MLA）采用另一种约束：把键和值共同生成自一个较低维的潜在向量。省略位置部分时，其内容分支可以写为
\begin{equation}
 \bm c_j=\bm W_D\bm h_j,\qquad
 \bm k_j^{C}=\bm W_{UK}\bm c_j,\qquad
 \bm v_j^{C}=\bm W_{UV}\bm c_j,
 \label{eq:nlp-llm-mla}
\end{equation}
其中$\bm c_j\in\mathbb R^{d_c}$，$d_c$小于展开后的多头维度。这里仍按单个内容头书写：$\bm W_D\in\mathbb R^{d_c\times d}$，$\bm W_{UK}\in\mathbb R^{d_h\times d_c}$与$\bm W_{UV}\in\mathbb R^{d_v\times d_c}$是该头对应的行块；各头共享同一$\bm c_j$，使用各自的上投影行块。对内容打分有$\bm q_t^{\top}\bm W_{UK}\bm c_j=(\bm W_{UK}^{\top}\bm q_t)^{\top}\bm c_j$，因此可以先变换当前查询，再在低维状态上完成匹配，不必为每个历史位置保存展开后的各头键。值聚合也可以利用线性性：
\begin{equation}
 \sum_{j=1}^{t}a_{tj}\bm W_{UV}\bm c_j
 =\bm W_{UV}\sum_{j=1}^{t}a_{tj}\bm c_j.
 \label{eq:nlp-llm-mla-value}
\end{equation}
这说明可以先在潜在向量上加权，再执行该头的值上投影，历史缓存同样不必保存逐头展开值。

位置旋转一般不能任意穿过投影矩阵。DeepSeek-V2因此把承载RoPE的位置查询与共享位置键从内容分支分开，分别拼接后计算打分；历史位置需要保存内容潜在向量和解耦的位置键，而不只是$\bm c_j$\scite{nlp-llm-deepseekv2-2024}。例如，在一个仅用于说明状态量的配置中，若$H=32$、$d_h=128$、$G=8$、$d_c=512$、位置键维度为64，则MHA、GQA与MLA分别保存8192、2048与576个数。这个比较只计算一层、一个历史词元的状态元素，不能据此直接推出整体运行时间的相同比例变化。

上述方案都允许查询读取全部历史位置，只改变匹配与状态的参数化。Transformer的通用结构可参考\BookChapterRef{chap:transformer}{卷五的Transformer章}；本节保留当前语言建模所需的计算，并把减少读取位置和汇入有限状态两种改变放到后面分别讨论。

\subsection{历史信息的选择性读取}
\label{sec:nlp-llm-sparse-reading}

若上下文长度为$T$，全注意力要为每个位置与其前缀计算匹配，匹配数量随$T^2$增长。键值压缩减少了每个历史位置的状态，却没有直接减少这些位置对。\term{稀疏注意力}（sparse attention）改为规定每个查询读取的历史子集或历史摘要，因而同时改变计算量与可用信息。最简单的滑动窗口只读取最近$w$个位置，匹配数量约为$Tw$；但若改期信息在窗口之外，模型就可能无法直接取到它。

压缩摘要与内容选择回应两种不同需要。把若干连续位置压成一个表示，可以用较少单元覆盖较远的历史，但具体细节可能在压缩时丢失。根据当前查询选择少量原始块，则保留被选块中的细节，但需要额外规则判断哪些块重要。二者不能都被描述成“只保留重要词元”：前者构造新表示，后者选择原位置。

原生稀疏注意力（Native Sparse Attention, NSA）把压缩、选择和局部窗口作为三条分别计算注意力的分支。设分支$b\in\{c,s,w\}$在当前位置得到输出$\bm o_t^{b}$，再用当前表示产生的门值$g_t^{b}\in[0,1]$融合：
\begin{equation}
 \bm o_t=g_t^{c}\bm o_t^{c}
             +g_t^{s}\bm o_t^{s}
             +g_t^{w}\bm o_t^{w}.
 \label{eq:nlp-llm-nsa}
\end{equation}
压缩分支提供粗粒度的全局信息和选块依据，选择分支读取得分较高块中的原始细节，局部分支负责邻近依赖。各分支分别归一化，不是把三类键值简单拼接后共用一次softmax；原文还为分支提供独立键值以减轻职责干扰\scite{nlp-llm-nsa2025}。

例如，一段很长的活动记录中，当前问题询问最终地点，选择分支有机会保留含地点更改的旧块，局部分支则读取最近的问句。选块若遗漏了关键旧块，局部计算再准确也无法恢复那条证据。块的离散选择还意味着，能够在预训练中学习整个方案，不等于每个选块决策都是普通的连续可微映射。讨论收益时应分别检查读取数量、选择开销与信息遗漏，而不能把理论稀疏性直接等同于特定设备上的加速。

\subsection{历史信息的递归汇聚}
\label{sec:nlp-llm-recurrent-state}

稀疏注意力仍保留可供选择的历史位置。另一种做法是不逐项保留全部历史，而是在读入每个词元时更新有限大小的状态。设$\bm s_t$是读到位置$t$后的状态，简化的\term{状态空间模型}（state space model, SSM）可写为
\begin{equation}
 \bm s_t=\bm A_t\bm s_{t-1}+\bm B_t\bm h_t,
 \qquad
 \bm o_t=\bm C_t\bm s_t.
 \label{eq:nlp-llm-state-space}
\end{equation}
$\bm A_t$控制旧信息怎样保留，$\bm B_t$控制当前表示怎样写入，$\bm C_t$控制读出。若状态为$m$维、输入为$d$维、输出为$d_o$维，则三者形状依次为$m\times m$、$m\times d$与$d_o\times m$。固定维度状态的大小不随已读长度增长；代价是许多历史细节需要共用同一状态，模型不能像全注意力那样直接访问每个旧位置。

Mamba让离散化步长及部分写入、读出参数依赖当前输入，使模型能够依据内容控制保留与遗忘。式\eqref{eq:nlp-llm-state-space}只概括状态更新关系，实际结构还包含受约束的状态矩阵、局部卷积与门控等部件。Mamba-2从结构化状态空间对偶性（structured state space duality, SSD）解释某些状态模型与结构化线性注意力的联系\scite{gu2023mamba,dao2024mamba2}。这个联系有指定的矩阵结构条件，不意味着所有状态空间模型等价于任意softmax注意力。

对需要按键修改记忆的情形，只进行加法累积还可能让旧值与新值互相干扰。令$\bm S_t\in\mathbb R^{d_v\times d_k}$表示键到值的关联状态，$\bm k_t\in\mathbb R^{d_k}$、$\bm v_t\in\mathbb R^{d_v}$是当前键值。先遗忘得到$\bar{\bm S}_t=\alpha_t\bm S_{t-1}$，再按当前键上的预测残差更新：
\begin{equation}
 \bm S_t=\bar{\bm S}_t+
 \beta_t(\bm v_t-\bar{\bm S}_t\bm k_t)\bm k_t^{\top},
 \qquad \bm o_t=\bm S_t\bm q_t.
 \label{eq:nlp-llm-gated-delta}
\end{equation}
其中$\alpha_t\in(0,1)$控制遗忘，$\beta_t\in(0,1)$控制修正强度，$\bm q_t$为读出查询。括号中的残差说明当前状态对这个键还欠缺多少值信息；更新只沿$\bm k_t$的方向修正。该式为便于理解采用标量遗忘门，实际变体可以使用更细的门控。

Gated DeltaNet将遗忘控制与这种增量修正规则结合；Kimi Linear进一步使用细粒度门控的状态更新，并与MLA层混合\scite{nlp-llm-gateddelta2024,nlp-llm-kimilinear2025}。混合结构能够让某些层汇聚历史、另一些层访问保留的位置，但其状态与缓存不能再统一称为固定大小。有限状态的优势和精确回读的需要之间存在取舍：若任务要求从很长前缀中精确复制任意旧字符串，状态容量与写入干扰就是必须检查的条件。

\subsection{逐位置变换与专家调用}
\label{sec:nlp-llm-experts}

上下文聚合回答当前位置从其他位置取得哪些信息；取得这些信息后，还需要对其进行非线性变换。\term{逐位置前馈网络}（position-wise feed-forward network, FFN）对每个位置使用相同参数，但在一次调用中不再让位置彼此交互。以门控变换为例，令输入为$\bm u_t\in\mathbb R^d$，可写为
\begin{equation}
 f(\bm u_t)=\bm W_2\bigl[
 \operatorname{SiLU}(\bm W_1\bm u_t)
 \odot(\bm W_3\bm u_t)\bigr],
 \label{eq:nlp-llm-ffn}
\end{equation}
其中$\bm W_1,\bm W_3\in\mathbb R^{m\times d}$，$\bm W_2\in\mathbb R^{d\times m}$，$\operatorname{SiLU}(z)=z\sigma(z)$。一条投影分支调节另一条分支的输出，最后回到$d$维。这是SwiGLU型变换；具体网络也可以使用其他激活与偏置约定\scite{shazeer2020glu}。

稠密前馈每次都调用整套参数。若希望提供更多可选择的变换，又不在每个词元上计算所有变换，可以准备$E$个专家$f_1,\ldots,f_E$，再由路由器根据$\bm u_t$选择少量专家。\term{混合专家}（mixture of experts, MoE）在这里表示按输入有条件地组合若干参数子网络，例如
\begin{equation}
 f_{\mathrm{MoE}}(\bm u_t)=
 \sum_{e\in\mathcal T_t}r_{t,e} f_e(\bm u_t),
 \label{eq:nlp-llm-moe}
\end{equation}
$\mathcal T_t$是路由得分最高的$k$个专家，$r_{t,e}$是它们的组合权重。总参数包含全部专家，某个词元的激活参数只包含实际调用的专家及共享部件。它们回答不同的资源问题，不应把总参数与稠密模型的每词元计算量直接比较\scite{shazeer2017}。

DeepSeekMoE将专家划分得更细，使固定激活量下有更多组合，并保留总会调用的共享专家来承接各输入共同需要的变换。共享专家之外，路由专家仍由输入选择\scite{nlp-llm-deepseekmoe2024}。这种设计提供分工机会，但“专家”只是网络部件的名称，不保证每个专家对应一个可解释领域。若多数词元集中到少量专家，其余参数就难以得到充分更新，因此需要检查路由比例及其随训练的变化。

负载控制可以通过辅助损失或路由调整实现，不过过强地追求均匀调用也可能妨碍有用分工。DeepSeek-V3报告了在主路由目标外调节专家选择偏置的负载控制方案\scite{deepseek2024v3}。这里关注它怎样改变专家的调用机会，而不展开训练设备之间的分配。专家计算仍需读入上下文表示并完成非线性映射，下一小节的条件记忆则改变参数信息的取用方式。

\subsection{条件参数记忆的接入}
\label{sec:nlp-llm-conditional-memory}

像常见短语、固定搭配和局部词形这样的模式会反复出现。若每次都只靠多层变换重新构造这些局部关系，可以考虑另设一条路径：用已经读到的词元片段索引可学习参数，再根据当前上下文决定取出的内容是否适用。这里的记忆是模型参数的一部分，与使用时检索外部文档不同，也不等于直接保存训练原文。

Engram采用多个词元后缀片段及多个哈希头查表。若$g_{t,n}$为截至$t$的$n$元片段，$\varphi_{n,k}$为第$k$个哈希函数，$\bm E_{n,k}$为对应可学习表，则一次读取为
\begin{equation}
 \bm e_{t,n,k}=\bm E_{n,k}[\varphi_{n,k}(g_{t,n})],
 \qquad
 \bm e_t=\operatorname{concat}_{n,k}(\bm e_{t,n,k}).
 \label{eq:nlp-llm-engram}
\end{equation}
确定性索引让相同规范化片段取到相同的表项，多组哈希减少所有读取都发生相同碰撞的机会。表项仍通过学习更新，哈希映射并没有赋予它们真实事实的保证\scite{nlp-llm-engram2026}。

取出的片段向量不依赖更远上下文，可能受到多义性和哈希碰撞影响。Engram因此将它投影为键和值，用已聚合上下文的隐藏表示产生门控，再通过短的因果卷积细化，最后残差接入主网络。门控回答“这个局部记忆在当前句子里应当贡献多少”，而查表回答“取哪一组参数”；两者职责不能合并成一次相似度检索。

与式\eqref{eq:nlp-llm-moe}相比，专家方案选择要运行的子网络，条件查表选择要读入的参数向量。两者可以共存，也都有容量、调用量与学习充分性之间的限制。Engram属于近期结构探索，其具体参数分配和任务收益需要在对应设置下检验。建立这种区分后，我们再考虑上述分支怎样在多层网络中传递。

\subsection{层间连接与多层组合}
\label{sec:nlp-llm-layer-composition}

单层聚合与变换只能完成一次信息处理。叠加层后，较高层可以使用已经融合上下文的表示，但也需要避免每层都完全改写输入。常用的前置归一化残差块把更新写成“旧表示加新分支”：
\begin{align}
 \bm u_t^{(\ell)}&=\bm h_t^{(\ell)}+
  \operatorname{Attn}_t\bigl(\operatorname{Norm}(\bm H^{(\ell)})\bigr),
 \label{eq:nlp-llm-layer-attention}\\
 \bm h_t^{(\ell+1)}&=\bm u_t^{(\ell)}+
  f\bigl(\operatorname{Norm}(\bm u_t^{(\ell)})\bigr).
 \label{eq:nlp-llm-layer-ffn}
\end{align}
$\bm H^{(\ell)}$表示这一层全部可见位置的表示，$\operatorname{Attn}_t$返回位置$t$的聚合输出，$f$可取稠密前馈或专家组合。归一化控制进入分支的数值尺度，残差路径使原表示不必完全经过该分支才能传下去；分支仍可改变表示，不能把残差理解为每层实际都接近恒等变换。

若采用均方根归一化，对$\bm h\in\mathbb R^d$的一种写法为$\operatorname{Norm}(\bm h)=\bm g\odot\bm h/\sqrt{d^{-1}\sum_i h_i^2+\varepsilon}$，其中$\bm g$是可学习尺度，$\varepsilon>0$用于避免分母为零。它不减去均值，与标准层归一化不同。若一条残差更新为$\bm h'=\bm h+F(\bm h)$，其局部导数为$\bm I+\partial F/\partial\bm h$，这说明梯度存在不经过$F$的直接路径；多层导数仍相乘，因而不能据此宣布梯度永远不会放大或衰减。

Hyper-Connections将单条残差流扩展为多条流，并学习流之间与计算分支之间的混合。为理解这种变化，设$\bm X^{(\ell)}\in\mathbb R^{r\times d}$保存$r$条流，可用示意式
\begin{equation}
 \bm X^{(\ell+1)}=
 \bm M_\ell\bm X^{(\ell)}+
 \bm b_\ell F(\bm a_\ell^{\top}\bm X^{(\ell)})
 \label{eq:nlp-llm-hyperconnection}
\end{equation}
表示其职责：$\bm a_\ell,\bm b_\ell\in\mathbb R^r$，$\bm M_\ell\in\mathbb R^{r\times r}$。$\bm a_\ell$将多流读入分支，$\bm b_\ell$把分支输出写回多流，$\bm M_\ell$控制旁路混合。实际方案还规定这些系数怎样依赖输入。多流增加组合自由度，也使反复混合的传播尺度成为需要检查的问题\scite{nlp-llm-hyperconnections2024}。

流形约束超连接（manifold-constrained hyper-connections, mHC）进一步约束残差混合矩阵，使理想的$\bm M_\ell$为双随机矩阵，即元素非负、每行和每列均为1。每条输出流因此是输入流的凸组合，多个这样的矩阵相乘也保持双随机性。这一性质限制的是残差混合本身，完整网络还有分支计算和其他映射。原文通过有限步Sinkhorn行列归一化近似实现约束；实际有限迭代的矩阵不必精确满足两个和为1的条件，不能把理想性质无条件搬到数值实现上\scite{nlp-llm-mhc2025}。

\subsection{词元概率与序列概率}
\label{sec:nlp-llm-probabilities}

经过$L$层计算后，位置$t$得到表示$\bm h_t^{(L)}$。要预测下一个词元，需要给词表中每个候选分配归一化概率。令最终归一化表示为$\bm z_t$，输出矩阵为$\bm U\in\mathbb R^{V\times d}$，偏置为$\bm b\in\mathbb R^V$，则
\begin{equation}
 p_\theta(x_{t+1}=v\mid x_{1:t})=
 \frac{\exp((\bm U\bm z_t+\bm b)_v)}
 {\sum_{u\in\mathcal V}\exp((\bm U\bm z_t+\bm b)_u)}.
 \label{eq:nlp-llm-softmax}
\end{equation}
参数$\theta$包含嵌入、各层变换与输出映射。某些模型让输出矩阵与输入嵌入共享参数，也有模型省去输出偏置；这些选择不改变这里从有限得分到词表分布的归一化关系。

概率链式法则将条件概率连成序列概率。加入开始标记后，可写为
\begin{equation}
 p_\theta(x_{1:T})=\prod_{t=1}^{T}
 p_\theta(x_t\mid x_{<t}).
 \label{eq:nlp-llm-sequence-probability}
\end{equation}
这里$x_{<t}$包括用于起始的约定条件。若要比较完整可变长度文本，应把结束标记也作为被预测事件，否则这个乘积只描述指定长度前缀的概率。softmax使每个局部词元分布归一化；模型是否几乎必然在有限步结束，还需要另外检查生成规则与终止性质。

在通知例子中，$p_\theta(\text{“周五”}\mid\text{“读书会改到”})$只比较这个位置的候选；整句概率还包含后续“下午”“两点”等事件的条件概率。单步最可能的词元未必构成整段最可能的序列，这一差别将在\ref{chap:nlp-model-inference}章用于解释解码。当前我们已经建立可求导的概率计算，但尚未规定每个样本中哪些输出必须正确，下一节才构造学习目标。

\section{预训练目标构造}
\label{sec:nlp-llm-objectives}

概率计算本身没有规定应当学会什么。相同网络可以只预测下一词元，也可以根据缺口两侧恢复中间片段；相同文本也可以产生多种监督记录。构造目标时，应分别回答：预测时允许看到哪些信息，要预测哪个对象，以及哪些错误按多大权重计入损失。三项决定需要共同明确，才能判断一个结构是否泄露答案、一个损失是否对应使用时的问题。

\subsection{可见条件与信息范围}
\label{sec:nlp-llm-objective-visibility}

对完整教学序列“读书会／改到／周五／下午／两点／举行”，若要学习左到右续写，预测“周五”时只能看到前两个词元。若要学习填空，可以把“周五”遮住，并同时提供左右两侧。两者预测对象看似相同，条件分布却不同：$p_\theta(x_3\mid x_1,x_2)$与$p_\theta(x_3\mid x_1,x_2,x_4,x_5,x_6)$不是同一个量。条件增加后问题可能更容易，但其损失也不能直接作为前一问题的替代。

设$M_{ij}\in\{0,1\}$表示输入位置$i$能否读取位置$j$。在注意力分数中，对$M_{ij}=0$的位置加上$-\infty$，再对允许位置归一化，就实现了可见范围。因果掩码取$M_{ij}=1$当且仅当$j\leq i$，位置$i$随后预测$x_{i+1}$；前缀方案可以让一段已给条件相互读取，同时限制待生成部分只能读取条件和自身过去。掩码的作用是限制信息流，而不是标记哪些位置计分，两者要分别设置。

另一种改变条件的方法是重排文本。例如先写出前缀与后缀，最后才放被移走的中间片段，那么普通因果计算在预测中间片段时，已经能够读取它两侧的内容。这里没有让网络读取计算序列中的未来，而是改变了原文位置与计算顺序的对应关系。训练记录必须保留特殊边界标记，否则模型无法判断当前是在续写原文还是补写缺口。

泄漏还可能发生在多层传播中。即使某个输出位置没有直接读取目标位置，若它先读取一个已经包含目标信息的中间表示，仍可能间接得到答案。因此，信息范围应针对整个计算图检查，不能只观察最后一层的掩码。需要排除的是目标通过本应不可见的位置或计算路径泄漏，而不是禁止一切复制。
例如从可见材料提取实体，或预测前文已经出现过的词元，都可以是合法的条件预测；
关键在于预测时使用的信息是否符合当前目标约定。

\subsection{预测对象与监督位置}
\label{sec:nlp-llm-objective-positions}

在可见条件确定后，还需要说明输出对应哪些目标。下一词元预测把输入位置$t$的表示对应到目标$x_{t+1}$；被遮蔽片段恢复把缺失位置对应到原始内容；多个未来位置预测则让同一前缀产生不同步长的输出。目标的编号必须与可见范围一致，不能把输入$x_t$对应到不受遮蔽的目标$x_t$，再误以为得到了语言预测任务。

令$\mathcal S(x)$表示一个样本中参与计分的目标位置集合，$C_i(x)$表示预测目标$x_i$时允许使用的条件。样本组织与目标构造共同产生记录$\{(C_i(x),x_i):i\in\mathcal S(x)\}$。在下一词元目标中，可以令$\mathcal S(x)=\{1,\ldots,T\}$、$C_i(x)=x_{<i}$；在遮蔽目标中，$\mathcal S(x)$通常只包含被遮蔽位置，条件包括未遮蔽内容及其位置。开始、结束、填充等特殊符号是否计分，也需要在这个集合中明确。

例如，把两篇文档装入一个样本后，填充位置通常不应作为普通正文目标计分；若文档边界强制隔离，第二篇文档的第一个词元也不能依赖第一篇的尾部。对重叠切分的长文档，还可以只对新引入部分计分，让重叠前缀提供条件而不重复贡献全部损失。改变计分位置会改变每段材料的实际学习权重，即使输入字符串没有变化。

多个未来位置预测还要区分辅助目标与生成协议。让当前位置同时预测$x_{t+1}$和$x_{t+2}$，会使共享表示收到两项误差信号；但若两个输出头都只读取$x_{1:t}$，第二个头就没有以第一个头实际选出的词元为条件。它们是给定同一前缀的两个预测分布，不能仅因共同训练就称为完整的未来联合分布。相应结构将在\ref{sec:nlp-llm-objective-comparison}小节写出。

\subsection{概率损失与目标组合}
\label{sec:nlp-llm-objective-loss}

现在可以用目标词元得到的概率为每项预测计分。若一次预测给真实目标分配的概率为$p$，负对数损失$-\log p$在$p$接近1时较小，在$p$接近0时增大。对一个样本及其可能的随机遮蔽，通用的计分形式为
\begin{equation}
 \ell(\theta;x)=
 -\frac{1}{Z(x)}\sum_{i\in\mathcal S(x)}
 w_i(x)\log p_\theta(x_i\mid C_i(x)),
 \label{eq:nlp-llm-objective-general}
\end{equation}
其中$w_i(x)\geq0$为位置权重，$Z(x)>0$给出归一化约定。若所有目标位置等权且$Z(x)=|\mathcal S(x)|$，得到每样本的平均目标词元损失；若先汇总一批样本的所有目标词元再平均，长样本对批次目标的贡献更大。两种口径需要明确，不能把它们混称为相同的平均损失。

对实际采样分布$P_{\bm\alpha}$及目标构造中的随机性取期望，得到学习的总体目标$\mathcal L(\theta)=\mathbb E[\ell(\theta;x)]$。训练时以抽到的有限批次估计其梯度。对于条件分布$q(\cdot\mid C)$，等权交叉熵可分解为该分布自身的熵与$D_{\mathrm{KL}}(q\|p_\theta)$之和，因此在固定条件、充分表达且能够求得最优值时，目标鼓励模型拟合数据条件分布。现实中$q$未知，有限材料与有限优化使这个解释不能变成对任意新输入的保证。

若同时使用若干目标，需说明每一项计算什么，再设置组合系数。例如$\mathcal L=\mathcal L_{\mathrm{next}}+\lambda\mathcal L_{\mathrm{future}}$让未来预测作为辅助信号，$\lambda\geq0$决定其相对作用。改变$\lambda$不仅改变数值大小，也改变共享参数收到的更新方向；一项损失下降可能伴随另一项上升。只比较总目标值，还可能掩盖各部分之间的取舍。

损失值还受到词元化、可见条件、目标位置和归一化的影响。具有双侧条件的填空模型得到较低损失，不足以说明它在左到右生成时更准确。即使协议完全相同，平均预测损失也可能被常见局部搭配主导，而用户关心的是少数关键条件或精确推理。预训练目标提供可学习的代理，具体能力需要独立检查。

\subsection{代表目标的完整对照}
\label{sec:nlp-llm-objective-comparison}

先固定同一序列$x_{1:6}$为上述通知短句。下一词元预测（next-token prediction, NTP）对各前缀预测下一个目标，损失可取$-\sum_{i=1}^{6}\log p_\theta(x_i\mid x_{<i})$，并按计分词元数平均。预测“周五”时看不到“下午两点”，预测“举行”时可以看到此前全部内容。所有目标共享一套概率映射，数据提供了大量不同长度的前缀。

中间填充（fill-in-the-middle, FIM）把原文分成前缀$A$、中间$B$和后缀$C$。若取$A=$“读书会改到”、$B=$“周五”、$C=$“下午两点举行”，可构造带边界标记的计算序列“前缀标记、$A$、后缀标记、$C$、中间标记、$B$”。在该序列上继续使用下一词元损失，预测$B$时就能读取$A$与$C$。FIM改变的是样本顺序与条件组织，仍可沿用因果结构；其混入比例与切分方式会影响学习结果\scite{nlp-llm-fim2022}。

UL2把多种去噪记录混合：常规跨度遮蔽只移走部分片段，顺序模式把后段作为待恢复内容，强遮蔽模式则隐藏更多内容或更长片段。它在编码器—解码器结构中，用模式标记说明当前恢复要求，再预测被移走内容\scite{nlp-llm-ul2-2022}。例如，输入“读书会改到、缺口、举行”，目标可为带片段标记的“周五下午两点”；换成顺序模式时，输入只保留前缀。这里改变了损坏方式与可见范围，不能把UL2理解为给普通NTP损失换了一个系数。

\term{多词元预测}（multi-token prediction, MTP）让同一前缀还承担更远未来的预测。共享主干、独立未来头的一种目标为
\begin{equation}
 \mathcal L_{\mathrm{MTP}}=
 -\sum_{j=1}^{J}\lambda_j
 \sum_{t:\,t+j\leq T}
 \log p_{\theta,j}(x_{t+j}\mid x_{1:t}),
 \label{eq:nlp-llm-mtp}
\end{equation}
其中$j$是预测步长，$J$是头数，$\lambda_j\geq0$为权重；这里省略了样本平均。以$x_{1:2}=$“读书会改到”为条件，两个头分别预测“周五”和“下午”。后一个头必须从同一前缀承担更远预测，因此对主干表示提出了额外要求，但两个头的边缘概率相乘并不自动给出真实未来联合分布\scite{nlp-llm-mtp2024}。

DeepSeek-V3采用顺序连接的MTP模块，将前一模块表示与后续词元嵌入共同用于更远预测，不应与独立未来头的条件完全混同\scite{deepseek2024v3}。辅助未来目标在部分代码、生成及算法任务上提供了收益证据，但收益受到模型规模、训练量和任务类型影响，并非增加预测头就必然改善所有任务。使用这些头提出多个候选词元，还需要\ref{chap:nlp-model-inference}章的验证与接受规则，训练目标本身不决定生成分布是否保持不变。

掩码扩散语言模型（masked diffusion language model）采用另一条生成路线：先规定怎样把完整文本逐步变成包含掩码的序列，再学习反向恢复。为说明一种常见设定，取噪声水平$\tau\sim\operatorname{Uniform}(0,1)$，将每个词元独立地以概率$\tau$替换为掩码符号，得到$\tilde x$。在线性遮蔽日程下，可用如下加权恢复目标：
\begin{equation}
 \mathcal L_{\mathrm{mask}}=
 \mathbb E_{\tau,x,\tilde x}
 \left[-\frac{1}{\tau}
 \sum_{i:\,\tilde x_i=\mathrm{mask}}
 \log p_\theta(x_i\mid\tilde x)\right].
 \label{eq:nlp-llm-masked-diffusion}
\end{equation}
期望覆盖完整文本、噪声水平及随机掩码。权重与所选前向损坏日程有关，不能任意移植到所有扩散模型；若使用不同日程、离散步数或预测参数化，就需要相应调整目标。

MDLM研究了掩码扩散的简化参数化与恢复损失，LLaDA在较大语言模型中采用掩码恢复路线，Dream则进一步探索扩散语言模型的学习与生成\scite{nlp-llm-mdlm2024,nlp-llm-llada2025,nlp-llm-dream2025}。对通知例子，受损输入可以同时保留“周五”和“举行”，要求恢复中间时间；当接近全掩码时，模型则必须更多依赖已学的文本分布。生成时通常从全掩码或部分已知序列出发，反复恢复与更新，网络输出的恢复分布并非式\eqref{eq:nlp-llm-sequence-probability}中的前缀条件分布，因此不能直接把这些输出按左到右顺序连乘来计算文本概率。完整采样过程留到\ref{chap:nlp-model-inference}章，当前的结论是：损坏过程、恢复目标和生成过程必须配套定义。

\section{预训练优化}
\label{sec:nlp-llm-optimization}

目标函数给出了希望降低的误差，尚未决定多少参数、多少训练词元能够在既定预算内配合，也未决定怎样把梯度变为参数变化。下面固定前面的训练分布与目标，依次讨论预算配置、参数尺度、更新方向与训练日程。这些属于学习算法及其条件；训练设备怎样通信和分配计算不在本节范围内。

\subsection{模型规模、词元规模与计算预算}
\label{sec:nlp-llm-compute-allocation}

在固定训练计算量下，扩大模型通常意味着能够处理的训练词元减少。若参数多而材料曝光不足，模型可能尚未充分学习；若模型太小，即使继续增加词元，也可能受表达能力限制。因此，比较规模时应同时给出参数量$N$、训练词元量$D$及计算预算$C$，而不是只用其中一个数字解释结果。

对结构与序列设置相近的稠密模型，可以用$C\approx\kappa ND$概括训练计算量，其中$\kappa$吸收一次前向与反向的系数。这是便于比较的近似，长上下文的注意力开销、输出映射和不同结构都可能使系数变化。Chinchilla在多种参数和词元规模上拟合训练结果，说明仅扩大参数而不相应增加训练词元，可能偏离所研究预算范围内较好的配置\scite{hoffmann2022chinchilla}。

一种便于理解的经验损失形式为
\begin{equation}
 L(N,D)=L_{\infty}+A N^{-p}+B D^{-q},
 \qquad A,B,p,q>0.
 \label{eq:nlp-llm-scaling-loss}
\end{equation}
$L_{\infty}$表示拟合中的不可约部分，后两项描述有限模型和有限数据带来的额外损失。代入$D=C/(\kappa N)$后，对$N$求驻点，得到
\begin{equation}
 pA N^{-p}=qB D^{-q},\qquad
 N_{\mathrm{opt}}\propto C^{q/(p+q)},\quad
 D_{\mathrm{opt}}\propto C^{p/(p+q)}.
 \label{eq:nlp-llm-scaling-allocation}
\end{equation}
第一式表示，在这个拟合模型和预算约束下，两种规模的边际改善需要平衡。若$p$与$q$接近，增加预算时两者大致同阶增长；这不是由概率论推导出的普遍定律，而依赖损失拟合、数据分布和可比较的训练过程。

训练词元量还不等于独立新材料量。把十亿词元重复十遍虽产生一百亿次曝光，但后几遍仍是相同材料，不能无条件套用不断加入新数据时的拟合关系。若同一模型将在训练后被调用很多次，总预算还包含使用计算，可能愿意多训练较小模型以降低每次使用成本。稀疏专家模型又需要分别记录总参数和激活参数；用哪一个量代入经验式，必须重新建立相应比较，不能直接继承稠密模型系数。

\subsection{参数化、初始化与更新尺度}
\label{sec:nlp-llm-parameter-scale}

选定规模后，相同的学习率数值未必在不同宽度上产生相同大小的函数变化。对一层的简化初始化分析，设输入$\bm h$为$d$维，一行权重$\bm w$的元素零均值、互相独立、方差为$s^2$，并独立于输入。给定$\bm h$时，线性输出满足$\operatorname{Var}_{\bm w}(\bm w^{\mathsf T}\bm h\mid\bm h)=s^2\lVert\bm h\rVert_2^2$。若各输入分量的二阶矩均为$\nu$，再对输入平均后，输出方差就是$ds^2\nu$；这是上述初始化假设下的计算，不表示训练后权重和激活仍保持独立。若宽度增大却保持$s^2$不变，初始激活尺度就会变化；让$s^2$随$d$缩小可以控制初始尺度，但训练更新还需另行检查。

一次参数变化$\Delta\theta$对输出的局部影响近似为$\Delta f\approx J_\theta\Delta\theta$，其中$J_\theta$是输出对参数的雅可比矩阵。即使$\Delta\theta$中每个元素大小相同，宽度变化也会改变参与相加的项数及$J_\theta$的尺度。因此，初始化合适并不自动保证更新合适；嵌入、隐藏矩阵、输出映射的形状不同，也不能把同一缩放规则不加区分地应用到所有参数。

最大更新参数化（maximal update parametrization, $\mu$P）协调不同参数组的初始化与学习率尺度，使宽度增加时隐藏特征的更新保持在相应的尺度上。$\mu$Transfer在这类参数化条件下，探索把较小模型上选择的超参数迁移到较大宽度，减少每个规模都完整调参的需要\scite{nlp-llm-mutransfer2022}。它处理的是跨规模的参数与更新关系，不是式\eqref{eq:nlp-llm-scaling-allocation}中模型与数据预算的最优分配。

迁移仍需固定哪些结构与目标保持一致，并在目标规模检查激活、梯度和损失。若改变层深、专家结构、词表、批次或训练分布，原有宽度关系未必充分描述新设置。可迁移的超参数也不意味着不需要验证；它提供的是有条件的调参组织方式，而非对任意模型免除优化检查的保证。

\subsection{梯度方向与矩阵优化}
\label{sec:nlp-llm-optimizer}

将全部参数按坐标排列后，每批样本给出梯度$\bm g^{(s)}$，优化器需要把它变成第$s$次更新。采用解耦权重衰减的Adam（Adam with decoupled weight decay, AdamW）分别维护一阶与二阶逐坐标统计，用历史梯度的尺度调整当前步长。省略偏差校正的计算细节后，其核心形式为
\begin{align}
 \bm m^{(s)}&=\beta_1\bm m^{(s-1)}+(1-\beta_1)\bm g^{(s)},\\
 \bm v^{(s)}&=\beta_2\bm v^{(s-1)}+(1-\beta_2)(\bm g^{(s)})^{\odot2},\\
 \theta^{(s+1)}&=(1-\eta_s\lambda)\theta^{(s)}
  -\eta_s\frac{\widehat{\bm m}^{(s)}}
  {\sqrt{\widehat{\bm v}^{(s)}}+\varepsilon}.
 \label{eq:nlp-llm-adamw}
\end{align}
其中$\beta_1,\beta_2\in[0,1)$为平滑系数，帽号表示从零初始化统计量时的偏差校正，即$\widehat{\bm m}^{(s)}=\bm m^{(s)}/(1-\beta_1^s)$、$\widehat{\bm v}^{(s)}=\bm v^{(s)}/(1-\beta_2^s)$。除法和平方根逐坐标进行，$\eta_s$是学习率。$\lambda\geq0$控制与梯度预处理分开的权重衰减，$\varepsilon>0$避免零分母。哪些参数使用衰减，应在参数组中说明\scite{loshchilov2019adamw}。

矩阵权重除了包含许多坐标，还定义输入到输出的线性方向。Muon对选定矩阵参数的动量方向做近似正交化，再调整更新尺度。为解释这一步，设矩阵动量的紧奇异值分解为$\bm M=\bm U\bm\Sigma\bm V^{\top}$，可以用$\bm Q=\bm U\bm V^{\top}$代替原方向。奇异向量保留，非零奇异值的大小差别被移除；实际实现通常以矩阵迭代近似这一变换，不需要逐次求完整分解\scite{jordan2024muon}。

这种变换与AdamW的逐坐标重标度不同：前者根据矩阵的输入输出方向处理更新，后者根据各坐标的历史平方梯度处理更新。正交化也不意味着新权重矩阵本身被约束为正交矩阵，因为被处理的是更新方向，更新仍加到原有权重上。对于非方阵，还需要说明所取尺度与形状的关系；同一谱方向若整体放大不同倍，参数变化量也不同。

《Muon is Scalable for LLM Training》研究了将Muon用于大规模语言模型时的权重衰减与更新尺度等调整，隐藏矩阵与嵌入、输出或其他非矩阵参数可使用不同更新规则\scite{nlp-llm-muonscale2025}。该工作并非最初提出Muon，也不能据其特定训练比较断言它对所有结构和目标均优于AdamW。合理比较需要固定数据、预算、调参范围和评价口径；优化器名称本身不能代替这些条件。

\subsection{批次、学习率与训练日程}
\label{sec:nlp-llm-training-schedule}

单次梯度依赖抽到哪些样本。若每个目标词元都等权，一批含$B$个计分词元时，梯度可写为它们负对数损失梯度的平均；若实际采用按序列平均，则批次定义随之改变。在近似独立、方差有限的理想情况下，增加批次能降低平均梯度的抽样方差，但长文档相邻位置和重复样本并不独立，计分词元数不能直接当作相同数量的独立观测。

学习率日程规定$\eta_s$怎样随学习进展变化。预热阶段从较小步长逐渐增加，可以限制训练初期统计尚未稳定时的参数变化；之后的衰减减小后期更新幅度。一种示例是先线性预热至$\eta_{\max}$，再按余弦曲线降到$\eta_{\min}$。选择这些曲线的目的在于组织探索与细调，不是证明任意网络会收敛到相同解；还应说明日程横轴采用更新次数还是已处理词元量。

假设两个配置都训练一百万次，但每次计分词元数相差四倍，它们并没有使用相同数据曝光或计算预算。若批次随训练增大，又仍按更新次数衰减，就可能改变每个学习率阶段处理的材料数量。记录时应把批次、累计词元、学习率、梯度尺度与留出损失联系起来，才能区分是学习进展、数据变化还是异常更新造成曲线变化。

DeepSeek-V3等技术报告同时给出结构、主辅目标、优化超参数与多阶段日程，适合作为完整组合的实例\scite{deepseek2024v3}。单份报告中的成功组合不能分离每个选择的独立贡献；把其中一个超参数移到不同结构，需要重新检查。训练过程产出的参数、数据曝光和优化记录，正是下一节进行阶段检查与继续学习的依据。

\section{继续预训练与阶段检查}
\label{sec:nlp-llm-continuation}

训练预算用完，只说明过程到达了预定终点。是否继续学习、采用哪组参数以及下一阶段还缺什么，需要依据未参与更新的数据和明确的能力问题判断。这里讨论预训练阶段内部的检查与适应；怎样构造完整评价协议、分析统计不确定性和污染风险，将在\ref{chap:nlp-evaluation}章统一展开。

\subsection{训练进展与参数选择}
\label{sec:nlp-llm-checkpoints}

训练损失反映当前采样批次的拟合，留出损失反映模型在同一约定下对未用于更新材料的预测。二者应采用相同的词元化、可见范围和计分方式；若留出材料的来源配比不同，还需分别报告来源结果。否则，整体损失变化可能仅来自比较对象变化，而非模型对同一分布预测得更好。

参数检查点是在训练过程中保存的参数及必要状态。选择检查点时，可以观察总体留出损失，并检查关心的来源和能力。若总体损失下降但某种少数语言持续恶化，按一个总均值选择参数就会掩盖该语言的退步。若多次根据留出结果修改配比和超参数，这批留出数据已经参与选择，最终结论还需要未反复用于决策的测试材料。

曲线平台期有多种可能解释：模型接近当前目标下的可达状态，学习率过小，材料过度重复，或者目标难以充分提供所需能力的信号。单凭曲线变平不能区分这些原因。应结合数据曝光、更新范数、各来源损失及异常记录检查，再决定增加数据、改变日程还是停止。突然的损失尖峰也应先排查样本与更新，而不能直接称为能力突然下降。

能力检查需要与本阶段目标相符。预训练模型可能正确续写知识性文本，却在直接问句后继续写另一个问题，因为它尚未被要求稳定扮演助手。把这种行为问题与知识缺失区分后，才能判断是需要更多预训练材料，还是应该交给有监督微调。检查的作用是定位下一步应改变的环节，不是用一个总榜单代替训练诊断。

\subsection{新领域与新语言的继续学习}
\label{sec:nlp-llm-domain-continuation}

若模型已经具有一般语言规律，但专业术语或某种语言覆盖不足，可以从现有参数出发，改用相关材料继续优化预训练目标。这称为\term{继续预训练}（continued pre-training）：仍从文本本身构造预测信号，主要变化在训练分布、曝光比例和更新安排。它与下一章根据指令—回答示范学习行为不同，虽然两种过程都从已有参数开始。

设原有材料分布为$P_{\mathrm{old}}$，新材料分布为$P_{\mathrm{new}}$，可以用$P_\rho=(1-\rho)P_{\mathrm{old}}+\rho P_{\mathrm{new}}$控制两者混合，$\rho\in[0,1]$。增大$\rho$使同一预算内更多更新针对新材料，但不保证旧分布目标同时下降。为解释这个冲突，设两类损失在当前参数的梯度分别为$\bm g_{\mathrm{old}}$和$\bm g_{\mathrm{new}}$。只沿新材料更新$\Delta\theta=-\eta\bm g_{\mathrm{new}}$时，旧损失的一阶变化近似为
\begin{equation}
 \Delta\mathcal L_{\mathrm{old}}
 \approx-\eta\langle\bm g_{\mathrm{old}},\bm g_{\mathrm{new}}\rangle.
 \label{eq:nlp-llm-continuation-conflict}
\end{equation}
若内积为负，这一步可能增加旧损失。该式是局部近似，解释目标冲突为何可能发生，不提供整个训练过程的变化保证。

保留一定旧材料、缩小更新强度或限制继续学习预算，都是控制改变范围的方式，但其具体作用需要新旧数据共同检查。对材料很少的新语言，单纯反复复制同一语料可能提高训练拟合，却增加记忆风险。词元化也影响条件：若新文字被切成极多字节片段，相同最大长度能容纳的语义内容会减少；扩展词表则还需要学习新增嵌入与输出参数，不能只改编码表而沿用原接口假设。

继续学习还需要区分事实更新与能力退步。活动从周四改到周五，模型对旧安排降低概率可能是正确更新；但因此失去读取其他活动时间的能力才是另一类问题。比较时应固定问题所依据的时间和材料，明确期望保留的是旧事实、通用技能还是输入协议，避免把一切行为变化都称为遗忘。

\subsection{上下文长度的继续适应}
\label{sec:nlp-llm-context-extension}

网络在公式上能够接收更长序列，不意味着它已经学会使用更远的证据。若训练中距离最大为$L$，使用时出现$sL$的序列，模型同时面对新位置范围、新距离组合和更多干扰内容。这里需要把位置机制的改变与长序列学习一起考虑，而不能只提高一个最大长度配置。

以式\eqref{eq:nlp-llm-rope}为例，直接把位置$t$扩展到原来范围之外，会引入训练中没有出现过的旋转组合。位置插值（positional interpolation）用$t/s$替代$t$，把更长序列的位置压回原有范围。若原最大位置为$L$、目标扩展倍数为$s>1$，则新序列末端$sL$映射到$L$附近；相邻位置的间隔也从1缩为$1/s$，所以局部相对关系同时发生变化。这解释了为什么插值可以限制位置范围，却仍需要通过继续学习适应新的距离分布\scite{nlp-llm-pi2023}。

YaRN采用随频率变化的缩放，综合保留部分原频率与插值后的频率，并调整注意力分数尺度，以减少统一压缩对不同频段的影响\scite{nlp-llm-yarn2023}。它是位置处理与学习安排的组合，不能仅凭频率公式保证任意长度的正确检索。评价时需要区分训练过的长度、插值范围以及进一步外推的长度；超过继续学习范围的结果不能与范围内结果合并成一个“支持长度”。

长序列样本也必须真正包含跨距离的预测需要。把大量互不相关短文档串成一条很长记录，虽然增加输入长度，却未必要求模型使用远距离关系。可以检查在不同距离放置改期通知、随后询问最终安排时，模型是否使用了正确证据，以及加入相似但无关安排后怎样变化。局部预测损失、跨段检索和长材料综合各考察不同问题，单个针式检索例子不能代表所有长文本能力。

\subsection{原有能力保持与阶段交接}
\label{sec:nlp-llm-pretraining-handoff}

完成初始或继续预训练后，应把新参数放回相同协议下与旧参数比较。对目标领域，检查相同留出材料上的损失和所需任务结果；对原有范围，检查语言覆盖、基本知识、长短上下文和必要行为边界。比较既要看变化量，也要看变化后的水平。例如，一个错误率从80\%降到60\%的模型虽然下降二十个百分点，仍不能据此称为已满足低错误率要求。

检查还应覆盖词元与位置接口。模型参数、词元化规则、特殊标记、结束符号和上下文配置需要作为一组产物保存；把不匹配的聊天模板或词表接到参数上，会产生与模型学习能力不同的运行问题。数据来源、时间范围、混合比例和训练目标说明，则帮助后续阶段判断哪些能力已经接触过、哪些材料仍不足。

进入有监督微调并不要求预训练模型已经具有完整助手行为，但需要知道其起点和缺口。若主要问题是对新领域语言缺少基本概率建模，少量交互示范未必足够补齐；若内容知识已经存在而输出方式不稳定，示范学习才直接针对这个行为差异。\ref{chap:nlp-supervised-finetuning}章接住的是这个明确的起点，\ref{chap:nlp-reasoning}章则把推理所需材料、轨迹与反馈放回多个学习阶段综合讨论。


\section{本章小结}
\label{sec:nlp-llm-summary}

预训练不是把大量文本交给一个足够大的网络就自然完成的过程。来源、筛选、词元化和混合确定模型实际反复观察的事件；上下文聚合、逐位置变换、参数记忆和层间连接确定可计算的条件概率；可见范围、预测对象与计分位置共同确定学习目标。模型结构和目标必须相容，例如因果预测不能泄露目标之后的信息，掩码扩散也必须配套定义从受损序列恢复文本的过程。

结构变化的收益需要放回它改变的对象中理解。键值共享与低维状态压缩减少历史状态的表示量，稀疏读取减少每次访问的位置，递归状态把历史汇入固定大小的记忆；这些方案保留信息的方式不同。专家调用和条件查表增加可用参数的方式也不同。参数更多、激活参数更少或者缓存更小，均不能单独推出任务正确率更高或端到端运行更快。

在给定分布和目标之后，规模选择、参数尺度、更新方向和训练日程共同决定参数能够学到什么。继续预训练允许针对领域、语言或长度缺口再次学习，但需要同时检查旧能力与新收益。交给下一阶段的产物应当包含参数、词元与位置协议、训练范围及检查记录。\ref{chap:nlp-supervised-finetuning}章将从这些产物出发，用示范回答约束交互行为；固定参数后的生成过程则在\ref{chap:nlp-model-inference}章展开。

\paragraph{延展阅读}
DataComp-LM与FineWeb可用于比较语料筛选和去重怎样改变训练分布，Chinchilla工作可用于
追踪固定计算预算下模型规模与词元规模的取舍，MDLM则提供不同于自回归预测的离散扩散
目标。继续阅读时，宜把数据、计算结构、训练目标和评价协议分别列出，再判断论文中的收益
能够归因到哪一项改变\scite{nlp-llm-dclm2024,nlp-llm-fineweb2024,hoffmann2022chinchilla,nlp-llm-mdlm2024}
